% 项目申请内容工作稿；未删节迁移稿见 proposal-full.tex / full.tex。
% 本轮扩写后的篇幅核验见 tests/ccnu-eduagenttrust-template/20260919-scientific-questions/REPORT.md。
% 本稿图文密集，收紧上下间距以避免整图后移和短尾页；不缩小原图。
\setlength{\intextsep}{2bp}
\HuazhongGuide{本表参照以下提纲撰写，要求逻辑清晰，层次分明，内容翔实，排版规范。（不超过7000 字）}

\section{项目的立项依据。研究目的、意义和国内外概况（附主要参考文献）}
\subsection{研究背景与意义}
ReAct和Toolformer展示了语言模型的行动与工具调用能力\cite{yao2023react,schick2023toolformer}，教育辅助进入连续执行，风险可能早于答案发生。同一解答在教师演示中适切，在独立练习中可能超界，需结合教学理论、学习者角色与授权判断\cite{zawacki2019review,ouyang2021paradigms}。

Bastani等报告缺少保护的辅助可能伴随撤去工具后表现下降，Kestin等报告教学原则指导的辅导取得积极即时结果\cite{bastani2025learning,kestin2025tutoring}；因对象和设计不同，不能直接外推高校智能体效果。结合模型局限与学习者能动性研究\cite{yan2024learning}，检验教育约束与行为证据的共同作用、跨任务稳定性及干预价值，支持工具选型与治理。
见图~\ref{fig:edu-f01}。
\HuazhongFigure{figures/F01-background-v2.pdf}{教育情境与过程证据共同决定风险含义}{fig:edu-f01}


\subsection{国内外研究现状}
\textbf{（1）教育智能体研究从回答质量逐步转向教学行为与任务过程。}Tutor CoPilot通过向真人辅导者提供实时建议，研究人工智能如何支持教学决策；其随机试验及对话分析发现，获得辅助的辅导者更倾向于使用引导性提问、减少直接给出答案，学生在相应知识点的掌握上也获得改善\cite{wang2024tutorcopilot}。这类研究将评价对象从回答是否正确推进到帮助方式是否适切，但其结论建立在人类辅导者参与的特定教学场景中，不能直接外推为自主工具型智能体的可靠性。

TeachArena预印本进一步将教学专业判断、多轮辅导和学习管理系统中的完整工作流程纳入共同评测，以逐轮回应、辅导轨迹及环境状态核验任务表现\cite{teacharena2026}。教育过程因此具有了更具体的评价对象：智能体既要理解学习者需要，也要在教学要求下完成行动。不过，教学行为评价、任务完成与学习效果属于不同终点；在这些能力评测基础上，仍需单独考察行为是否超出本次允许的帮助程度，以及风险证据能否在动作提交前被识别。

\textbf{（2）通用智能体评测已覆盖交互与安全风险，教育约束的作用仍需辨明。}清华大学等团队提出的AgentBench通过多种交互环境考察推理、决策和指令遵循，WebArena则在可复现网页环境中评价多步骤任务的功能正确性\cite{liu2023agentbench,zhou2023webarena}。二者为观察连续行为和环境后果提供了方法基础。ToolEmu利用语言模型模拟工具执行并评估潜在失败，降低复杂风险场景的测试成本；R-Judge依据带有安全标签和风险描述的多轮交互记录，直接考察模型识别行为风险的能力\cite{ruan2023toolemu,yuan2024rjudge}。由此可见，轨迹评价已有明确研究基础，教育场景需要进一步说明的是教学目的、角色权限与行为证据之间的关系。

在攻击与防御方面，AgentDojo和InjecAgent重点考察工具返回内容中的间接提示注入，AgentHarm则评测明确恶意请求下的多步骤有害行为\cite{debenedetti2024agentdojo,zhan2024injecagent,andriushchenko2024agentharm}。这些基准揭示了外部信息和恶意指令造成的行为偏移，但其威胁模型与正常学习请求中的支持失配有所不同。教育智能体即使未受到攻击，也可能过度代答、误用证据或越权访问；识别这些行为需要同时读取教学要求与执行记录，恶意请求上的拒绝表现不能替代这种情境判断。

\textbf{（3）过程监督与自动评判提供了评价手段，其有效性需要按对象验证。}逐步监督研究通过对中间推理步骤提供反馈改善数学推理，ProcessBench进一步将首个错误步骤的定位作为评测任务\cite{lightman2023verify,processbench2024}。这些工作说明中间证据具有独立评价价值，但推理错误与教育风险并不等同：某项操作可能逻辑正确、执行成功，却违反帮助边界或数据授权。过程评价因此还需回答“依据哪项约束、何处证据构成违背”，而不只是寻找错误步骤。

大模型自动评判研究同时揭示了评价器本身的局限。MT-Bench相关工作发现位置、冗长和自我偏好等偏差，G-Eval在改善与人工判断一致性的同时，也提示对模型生成文本的潜在偏好\cite{zheng2023judging,liu2023geval}。FActScore核查原子事实是否获得来源支持，Ragas区分检索相关性、证据利用与生成质量，TruthfulQA考察模型是否复述错误信念\cite{min2023factscore,es2023ragas,lin2021truthfulqa}。这些方法可为证据真实性和内容质量提供局部判据；将其用于教育可信评价时，还须验证与教学要求的对应关系，并区分事实错误、支持失配及权限违背，避免一个综合分数掩盖不同风险。

\textbf{（4）运行保护已发展出规则与分级控制，干预收益仍需结合任务效用解释。}GuardAgent将安全要求转化为检查计划和可执行保护代码，用于核验目标智能体行为；NEXUS预印本结合结构化计划、确定性规则、参数检查和校准风险评分，选择放行、阻断、请求确认或请求修订等处置\cite{xiang2024guardagent,agentguard2026nexus}。这些研究表明，动作检查和分级处置已具有可借鉴的技术路径。面向教育任务，进一步的问题是能否在证据充分且动作尚未提交时采取适当干预，并保留正常任务所需的帮助。

安全基准测量效度审计预印本指出，不同基准对应的目标行为并不一致，类别比例、任务能力与模型样本构成都会影响安全分数的解释\cite{wang2026safetyvalidity}。据此，较低的风险执行率还需与正常完成率、误阻断和人工接管负担共同观察；仅因更多拒绝而减少风险，尚不足以说明干预更有效。总体来看，现有研究分别提供了教学行为评测、轨迹风险判断和运行保护方法，而教育约束如何改变证据含义、有效证据何时可用、识别结果能否带来受效用约束的干预收益，仍需要相互衔接的验证。

\subsection{研究不足与科学假说}
现有研究已涉及教学行为、轨迹风险与运行保护，但教学支持边界对风险判定的作用仍需解释。同一解答在教师演示中适切，在独立练习中可能构成过度代答；答案正确性难以反映这种差异。完整轨迹中的风险未必能在动作前识别，风险减少还可能以额外拒绝和任务中断为代价。

关键科学问题是：教学约束如何改变行为证据的风险含义，何种条件下能支持兼顾任务效用的动作前干预？

本项目提出科学假说：在模型和预算一致时，正确教学约束能够增强轨迹证据对支持边界风险的识别贡献，该增益大于约束打乱时的相应增益；当有效证据先于动作提交出现时，分级干预可在正常任务自主完成率损失受限的条件下降低风险执行率。

通过交叉消融、统一名义误报校准、家族留出及干预对照，排查风险词、模板及额外拒绝的影响。依据效应区间及实际意义幅度判定支持程度，区间跨界保留不确定性，并由家族差异界定适用范围。
% 由 myexample.bib 和仓库共享 gbt7714-nsfc.bst 生成；28条被引文献。
% 原第9条 learnlm2024 已按用户要求删除；其余条目保持原相对顺序。
% 重生成：python3 tests/ccnu-eduagenttrust-template/20260919-references/render_references.py
\begin{thebibliography}{10}
\providecommand{\bibauthor}[1]{#1}
\providecommand{\bibeditor}[1]{#1}
\providecommand{\bibtranslator}[1]{#1}
\providecommand{\bibtitle}[1]{#1}
\providecommand{\bibbooktitle}[1]{#1}
\providecommand{\bibjournal}[1]{#1}
\providecommand{\bibmark}[1]{#1}
\providecommand{\bibcountry}[1]{#1}
\providecommand{\bibpatentid}[1]{#1}
\providecommand{\bibedition}[1]{#1}
\providecommand{\biborganization}[1]{#1}
\providecommand{\bibaddress}[1]{#1}
\providecommand{\bibpublisher}[1]{#1}
\providecommand{\bibinstitution}[1]{#1}
\providecommand{\bibschool}[1]{#1}
\providecommand{\bibvolume}[1]{#1}
\providecommand{\bibnumber}[1]{#1}
\providecommand{\bibversion}[1]{#1}
\providecommand{\bibpages}[1]{#1}
\providecommand{\bibmodifydate}[1]{#1}
\providecommand{\bibcitedate}[1]{#1}
\providecommand{\bibyear}[1]{#1}
\providecommand{\bibdate}[1]{#1}
\providecommand{\biburl}[1]{\newline\url{#1}}

\bibitem{yao2023react}
\bibauthor{Yao S\@, Zhao J\@, Yu D\@, Du N\@, Shafran I\@, Narasimhan K\@, Cao
  Y}\@. \bibtitle{{ReAct: Synergizing Reasoning and Acting in Language
  Models}}\bibmark{[J/OL]}\@. \bibjournal{arXiv preprint arXiv:2210.03629}\@,
  \bibyear{2022}\bibcitedate{[2026-09-06]}\@.
  \url{https://arxiv.org/abs/2210.03629}\@.

\bibitem{schick2023toolformer}
\bibauthor{Schick T\@, Dwivedi-Yu J\@, Dessì R\@, Raileanu R\@, Lomeli M\@,
  Zettlemoyer L\@, Cancedda N\@, Scialom T}\@. \bibtitle{{Toolformer: Language
  Models Can Teach Themselves to Use Tools}}\bibmark{[J/OL]}\@.
  \bibjournal{arXiv preprint arXiv:2302.04761}\@,
  \bibyear{2023}\bibcitedate{[2026-09-06]}\@.
  \url{https://arxiv.org/abs/2302.04761}\@.

\bibitem{zawacki2019review}
\bibauthor{Zawacki-Richter O\@, Marín V~I\@, Bond M\@, Gouverneur F}\@.
  \bibtitle{{Systematic review of research on artificial intelligence
  applications in higher education -- where are the
  educators?}}\bibmark{[J/OL]}\@. \bibjournal{International Journal of
  Educational Technology in Higher Education}\@,
  \bibvolume{16}\bibnumber{(1)}\thinspace{}\textnormal{: }\bibpages{39}\@,
  \bibyear{2019}\bibcitedate{[2026-09-06]}\@.
  \url{https://doi.org/10.1186/s41239-019-0171-0}\@.

\bibitem{ouyang2021paradigms}
\bibauthor{Ouyang F\@, Jiao P}\@. \bibtitle{{Artificial intelligence in
  education: The three paradigms}}\bibmark{[J/OL]}\@. \bibjournal{Computers and
  Education: Artificial Intelligence}\@, \bibvolume{2}\thinspace{}\textnormal{:
  }\bibpages{100020}\@, \bibyear{2021}\bibcitedate{[2026-09-06]}\@.
  \url{https://doi.org/10.1016/j.caeai.2021.100020}\@.

\bibitem{bastani2025learning}
\bibauthor{Bastani H\@, Bastani O\@, Sungu A\@, Ge H\@, Özge Kabakcı\@,
  Mariman R}\@. \bibtitle{{Generative AI without guardrails can harm learning:
  Evidence from high school mathematics}}\bibmark{[J/OL]}\@.
  \bibjournal{Proceedings of the National Academy of Sciences}\@,
  \bibvolume{122}\bibnumber{(26)}\thinspace{}\textnormal{:
  }\bibpages{e2422633122}\@, \bibyear{2025}\bibcitedate{[2026-09-06]}\@.
  \url{https://doi.org/10.1073/pnas.2422633122}\@.

\bibitem{kestin2025tutoring}
\bibauthor{Kestin G\@, Miller K\@, Klales A\@, Milbourne T\@, Ponti G}\@.
  \bibtitle{{AI tutoring outperforms in-class active learning: an RCT
  introducing a novel research-based design in an authentic educational
  setting}}\bibmark{[J/OL]}\@. \bibjournal{Scientific Reports}\@,
  \bibvolume{15}\bibnumber{(1)}\thinspace{}\textnormal{: }\bibpages{17458}\@,
  \bibyear{2025}\bibcitedate{[2026-09-06]}\@.
  \url{https://doi.org/10.1038/s41598-025-97652-6}\@.

\bibitem{yan2024learning}
\bibauthor{Yan L\@, Greiff S\@, Teuber Z\@, Gašević D}\@. \bibtitle{{Promises
  and challenges of generative artificial intelligence for human
  learning}}\bibmark{[J/OL]}\@. \bibjournal{Nature Human Behaviour}\@,
  \bibvolume{8}\bibnumber{(10)}\thinspace{}\textnormal{:
  }\bibpages{1839\thinspace{}\textnormal{--}\thinspace{}1850}\@,
  \bibyear{2024}\bibcitedate{[2026-09-06]}\@.
  \url{https://doi.org/10.1038/s41562-024-02004-5}\@.

\bibitem{wang2024tutorcopilot}
\bibauthor{Wang R~E\@, Ribeiro A~T\@, Robinson C~D\@, Loeb S\@, Demszky D}\@.
  \bibtitle{{Tutor CoPilot: A Human-AI Approach for Scaling Real-Time
  Expertise}}\bibmark{[J/OL]}\@. \bibjournal{arXiv preprint
  arXiv:2410.03017}\@, \bibyear{2024}\bibcitedate{[2026-09-06]}\@.
  \url{https://arxiv.org/abs/2410.03017}\@.

\bibitem{teacharena2026}
\bibauthor{Chen Z\@, Liu P\@, Sheng R\@, Li H\@, Tu J\@, Deng X\@, Shum K\@,
  Liu D\@, Qu H}\@. \bibtitle{{TeachArena: Are Language Agents Ready for
  Realistic Teaching Work?}}\bibmark{[J/OL]}\@. \bibjournal{arXiv preprint
  arXiv:2605.14322}\@, \bibyear{2026}\bibcitedate{[2026-09-06]}\@.
  \url{https://arxiv.org/abs/2605.14322}\@.

\bibitem{liu2023agentbench}
\bibauthor{Liu X\@, Yu H\@, Zhang H\@, Xu Y\@, Lei X\@, Lai H\@, Gu Y\@, Ding
  H\@, Men K\@, Yang K\@, Zhang S\@, Deng X\@, Zeng A\@, Du Z\@, Zhang C\@,
  Shen S\@, Zhang T\@, Su Y\@, Sun H\@, Huang M\@, Dong Y\@, Tang J}\@.
  \bibtitle{{AgentBench: Evaluating LLMs as Agents}}\bibmark{[J/OL]}\@.
  \bibjournal{arXiv preprint arXiv:2308.03688}\@,
  \bibyear{2023}\bibcitedate{[2026-09-06]}\@.
  \url{https://arxiv.org/abs/2308.03688}\@.

\bibitem{zhou2023webarena}
\bibauthor{Zhou S\@, Xu F~F\@, Zhu H\@, Zhou X\@, Lo R\@, Sridhar A\@, Cheng
  X\@, Ou T\@, Bisk Y\@, Fried D\@, Alon U\@, Neubig G}\@. \bibtitle{{WebArena:
  A Realistic Web Environment for Building Autonomous
  Agents}}\bibmark{[J/OL]}\@. \bibjournal{arXiv preprint arXiv:2307.13854}\@,
  \bibyear{2023}\bibcitedate{[2026-09-06]}\@.
  \url{https://arxiv.org/abs/2307.13854}\@.

\bibitem{ruan2023toolemu}
\bibauthor{Ruan Y\@, Dong H\@, Wang A\@, Pitis S\@, Zhou Y\@, Ba J\@, Dubois
  Y\@, Maddison C~J\@, Hashimoto T}\@. \bibtitle{{Identifying the Risks of LM
  Agents with an LM-Emulated Sandbox}}\bibmark{[J/OL]}\@. \bibjournal{arXiv
  preprint arXiv:2309.15817}\@, \bibyear{2023}\bibcitedate{[2026-09-06]}\@.
  \url{https://arxiv.org/abs/2309.15817}\@.

\bibitem{yuan2024rjudge}
\bibauthor{Yuan T\@, He Z\@, Dong L\@, Wang Y\@, Zhao R\@, Xia T\@, Xu L\@,
  Zhou B\@, Li F\@, Zhang Z\@, Wang R\@, Liu G}\@. \bibtitle{{R-Judge:
  Benchmarking Safety Risk Awareness for LLM Agents}}\bibmark{[J/OL]}\@.
  \bibjournal{arXiv preprint arXiv:2401.10019}\@,
  \bibyear{2024}\bibcitedate{[2026-09-06]}\@.
  \url{https://arxiv.org/abs/2401.10019}\@.

\bibitem{debenedetti2024agentdojo}
\bibauthor{Debenedetti E\@, Zhang J\@, Balunović M\@, Beurer-Kellner L\@,
  Fischer M\@, Tramèr F}\@. \bibtitle{{AgentDojo: A Dynamic Environment to
  Evaluate Prompt Injection Attacks and Defenses for LLM
  Agents}}\bibmark{[J/OL]}\@. \bibjournal{arXiv preprint arXiv:2406.13352}\@,
  \bibyear{2024}\bibcitedate{[2026-09-06]}\@.
  \url{https://arxiv.org/abs/2406.13352}\@.

\bibitem{zhan2024injecagent}
\bibauthor{Zhan Q\@, Liang Z\@, Ying Z\@, Kang D}\@. \bibtitle{{I}njec{A}gent:
  Benchmarking Indirect Prompt Injections in Tool-Integrated Large Language
  Model Agents}\bibmark{[C/OL]} //\thinspace{}\bibeditor{Ku L-W\@, Martins A\@,
  Srikumar V}\@. \bibbooktitle{Findings of the Association for Computational
  Linguistics: ACL 2024}\@. \bibaddress{Bangkok,
  Thailand}\thinspace{}\textnormal{: }\bibpublisher{Association for
  Computational Linguistics}\@, \bibyear{2024}\thinspace{}\textnormal{:
  }\bibpages{10471\thinspace{}\textnormal{--}\thinspace{}10506}\bibcitedate{[2026-09-06]}\@.
  \url{https://aclanthology.org/2024.findings-acl.624/}\@.

\bibitem{andriushchenko2024agentharm}
\bibauthor{Andriushchenko M\@, Souly A\@, Dziemian M\@, Duenas D\@, Lin M\@,
  Wang J\@, Hendrycks D\@, Zou A\@, Kolter Z\@, Fredrikson M\@, Winsor E\@,
  Wynne J\@, Gal Y\@, Davies X}\@. \bibtitle{{AgentHarm: A Benchmark for
  Measuring Harmfulness of LLM Agents}}\bibmark{[J/OL]}\@. \bibjournal{arXiv
  preprint arXiv:2410.09024}\@, \bibyear{2024}\bibcitedate{[2026-09-06]}\@.
  \url{https://arxiv.org/abs/2410.09024}\@.

\bibitem{lightman2023verify}
\bibauthor{Lightman H\@, Kosaraju V\@, Burda Y\@, Edwards H\@, Baker B\@, Lee
  T\@, Leike J\@, Schulman J\@, Sutskever I\@, Cobbe K}\@. \bibtitle{{Let's
  Verify Step by Step}}\bibmark{[J/OL]}\@. \bibjournal{arXiv preprint
  arXiv:2305.20050}\@, \bibyear{2023}\bibcitedate{[2026-09-06]}\@.
  \url{https://arxiv.org/abs/2305.20050}\@.

\bibitem{processbench2024}
\bibauthor{Zheng C\@, Zhang Z\@, Zhang B\@, Lin R\@, Lu K\@, Yu B\@, Liu D\@,
  Zhou J\@, Lin J}\@. \bibtitle{{ProcessBench: Identifying Process Errors in
  Mathematical Reasoning}}\bibmark{[J/OL]}\@. \bibjournal{arXiv preprint
  arXiv:2412.06559}\@, \bibyear{2024}\bibcitedate{[2026-09-06]}\@.
  \url{https://arxiv.org/abs/2412.06559}\@.

\bibitem{zheng2023judging}
\bibauthor{Zheng L\@, Chiang W-L\@, Sheng Y\@, Zhuang S\@, Wu Z\@, Zhuang Y\@,
  Lin Z\@, Li Z\@, Li D\@, Xing E~P\@, Zhang H\@, Gonzalez J~E\@, Stoica I}\@.
  \bibtitle{{Judging LLM-as-a-Judge with MT-Bench and Chatbot
  Arena}}\bibmark{[J/OL]}\@. \bibjournal{arXiv preprint arXiv:2306.05685}\@,
  \bibyear{2023}\bibcitedate{[2026-09-06]}\@.
  \url{https://arxiv.org/abs/2306.05685}\@.

\bibitem{liu2023geval}
\bibauthor{Liu Y\@, Iter D\@, Xu Y\@, Wang S\@, Xu R\@, Zhu C}\@.
  \bibtitle{{G-Eval: NLG Evaluation using GPT-4 with Better Human
  Alignment}}\bibmark{[J/OL]}\@. \bibjournal{arXiv preprint
  arXiv:2303.16634}\@, \bibyear{2023}\bibcitedate{[2026-09-06]}\@.
  \url{https://arxiv.org/abs/2303.16634}\@.

\bibitem{min2023factscore}
\bibauthor{Min S\@, Krishna K\@, Lyu X\@, Lewis M\@, Yih W-t\@, Koh P~W\@,
  Iyyer M\@, Zettlemoyer L\@, Hajishirzi H}\@. \bibtitle{{FActScore:
  Fine-grained Atomic Evaluation of Factual Precision in Long Form Text
  Generation}}\bibmark{[J/OL]}\@. \bibjournal{arXiv preprint
  arXiv:2305.14251}\@, \bibyear{2023}\bibcitedate{[2026-09-06]}\@.
  \url{https://arxiv.org/abs/2305.14251}\@.

\bibitem{es2023ragas}
\bibauthor{Es S\@, James J\@, Espinosa-Anke L\@, Schockaert S}\@.
  \bibtitle{{Ragas: Automated Evaluation of Retrieval Augmented
  Generation}}\bibmark{[J/OL]}\@. \bibjournal{arXiv preprint
  arXiv:2309.15217}\@, \bibyear{2023}\bibcitedate{[2026-09-06]}\@.
  \url{https://arxiv.org/abs/2309.15217}\@.

\bibitem{lin2021truthfulqa}
\bibauthor{Lin S\@, Hilton J\@, Evans O}\@. \bibtitle{{TruthfulQA: Measuring
  How Models Mimic Human Falsehoods}}\bibmark{[J/OL]}\@. \bibjournal{arXiv
  preprint arXiv:2109.07958}\@, \bibyear{2021}\bibcitedate{[2026-09-06]}\@.
  \url{https://arxiv.org/abs/2109.07958}\@.

\bibitem{xiang2024guardagent}
\bibauthor{Xiang Z\@, Zheng L\@, Li Y\@, Hong J\@, Li Q\@, Xie H\@, Zhang J\@,
  Xiong Z\@, Xie C\@, Yang C\@, Song D\@, Li B}\@. \bibtitle{{GuardAgent:
  Safeguard LLM Agents by a Guard Agent via Knowledge-Enabled
  Reasoning}}\bibmark{[J/OL]}\@. \bibjournal{arXiv preprint
  arXiv:2406.09187}\@, \bibyear{2024}\bibcitedate{[2026-09-06]}\@.
  \url{https://arxiv.org/abs/2406.09187}\@.

\bibitem{agentguard2026nexus}
\bibauthor{Hossain E\@, Nipu M~M~H\@, Ornee T~N\@, Rana R\@, Yousefi N}\@.
  \bibtitle{{NEXUS: Structured Runtime Safety for Tool-Using LLM
  Agents}}\bibmark{[J/OL]}\@. \bibjournal{arXiv preprint arXiv:2607.19356}\@,
  \bibyear{2026}\bibcitedate{[2026-09-06]}\@.
  \url{https://arxiv.org/abs/2607.19356}\@.

\bibitem{wang2026safetyvalidity}
\bibauthor{Wang Y\@, Han X\@, Shang D\@, Tang Y\@, Liu B}\@. \bibtitle{{Safety,
  or Just Capability? A Validity Audit of Agent-Safety
  Benchmarks}}\bibmark{[J/OL]}\@. \bibjournal{arXiv preprint
  arXiv:2607.28685}\@, \bibyear{2026}\bibcitedate{[2026-09-06]}\@.
  \url{https://arxiv.org/abs/2607.28685}\@.

\bibitem{yan2024knowledge}
\bibauthor{Yan F\@, Li Z\@, Silamu W\@, Li Y}\@. \bibtitle{{Knowledge-aware
  image understanding with multi-level visual representation enhancement for
  visual question answering}}\bibmark{[J/OL]}\@. \bibjournal{Machine
  Learning}\@, \bibvolume{113}\bibnumber{(6)}\thinspace{}\textnormal{:
  }\bibpages{3789\thinspace{}\textnormal{--}\thinspace{}3805}\@,
  \bibyear{2024}\bibcitedate{[2026-09-06]}\@.
  \url{https://doi.org/10.1007/s10994-023-06426-6}\@.

\bibitem{tang2026hashing}
\bibauthor{Tang L\@, Wang L\@, Wang G\@, Sun M\@, Yan F}\@. \bibtitle{{Feature
  Fusion Mamba Hashing via Decoupling for Cross-Modal
  Retrieval}}\bibmark{[J/OL]}\@. \bibjournal{IEEE Signal Processing Letters}\@,
  \bibvolume{33}\thinspace{}\textnormal{:
  }\bibpages{609\thinspace{}\textnormal{--}\thinspace{}613}\@,
  \bibyear{2026}\bibcitedate{[2026-09-06]}\@.
  \url{https://doi.org/10.1109/LSP.2026.3652952}\@.

\end{thebibliography}


\section{项目的研究内容、研究目标,以及拟解决的关键科学问题。（此部分为重点阐述内容）}
\subsection{总体框架与风险表征}
围绕课程答疑和文献学习辅助，研究含检索、记忆组件的工具型单智能体；双智能体仅作传播压力测试。
见图~\ref{fig:edu-f02}。
\HuazhongFigure{figures/F02-framework-v2.pdf}{风险表征、轨迹评测与监测干预的总体框架}{fig:edu-f02}

R1按图中三类风险关联约束与证据，比较同答案异过程、同过程异授权的案例；由未参与生成的教学人员复核，保留争议。
见图~\ref{fig:edu-f03}。
\HuazhongFigure{figures/F03-risk-representation-v2.pdf}{R1：任务约束、风险证据与独立标注}{fig:edu-f03}


\subsection{过程级可信评测与动作前监测}
R2以外显计划、检索证据、工具参数、回执及记忆读写关联任务约束，不索取私有思维链。规则与语义评价的分工、证据分级见图。离线读取完整轨迹，在线只读当前前缀与候选动作；整轨迹取各维度已见最高分。正常校准轨迹用于选取不超过误报目标的最近可行阈值。
见图~\ref{fig:edu-f04}。
\HuazhongFigure{figures/F04-trajectory-evaluation-v2.pdf}{R2：当前前缀评测与动作前干预窗口}{fig:edu-f04}

R3承接R2的风险、证据与不确定性，按图中顺序检查候选动作；候选回答先缓存，发布即提交，无工具代答也须检查。校准时先筛选满足任务损失约束的策略，再比较风险执行率；无可行策略则报告失效。
见图~\ref{fig:edu-f05}。
\HuazhongFigure{figures/F05-monitoring-intervention-v2.pdf}{R3：候选动作缓存、顺序检查与分级处置}{fig:edu-f05}


\subsection{研究目标与关键科学问题}
本项目旨在阐明教学约束、行为证据与风险判断之间的关系，明确过程评价能够支持动作前干预的条件及其效用边界。围绕这一目标，建立可复核的风险案例集，形成情境约束与轨迹证据关联的评价方法，并通过留出任务检验干预判据。重点研究以下三个相互衔接的科学问题。

\textbf{第一，教学约束如何改变行为证据的风险含义？}同一回答或工具行动在不同帮助边界下可能具有不同的合规性，单靠答案正确性或风险词难以解释这种变化。需要辨明：正确的教学约束能否增强轨迹证据的判别作用，这种作用在任务表述和家族改变后能否保持。通过授权变化任务对与约束、证据交叉消融，区分情境与证据的共同贡献及提示长度、关键词等因素的影响。

\textbf{第二，完整轨迹中的可识别风险，在何种条件下能够提前判断？}事后能够确认的风险，未必在动作提交前已有充分证据。需要厘清证据出现、风险告警与动作提交之间的时序关系，判断哪些证据能够在当前前缀中支持告警，哪些风险只能事后识别。通过同一事件的前缀评价与完整轨迹评价，考察证据可用性、检测延迟及干预窗口对提前检出的限制，避免将事后信息带来的识别能力误作在线能力。

\textbf{第三，及时识别在何种条件下能够转化为有效干预？}提前告警并不必然降低风险，频繁拒绝或人工接管也可能以牺牲正常任务完成换取表面安全。需要揭示风险证据、可用干预窗口与处置方式如何共同影响运行结果。通过相同初始任务的策略分支比较，检验分级干预能否在正常任务自主完成率损失受限的条件下降低风险执行率，并界定其失效情形。

\section{拟采取的研究方案及可行性分析。（包括有关方法、技术路线、实验手段、关键技术等说明）}
\subsection{任务构建、标注与数据划分}
以许可明确的公开课程资料和论文为基础，围绕课程答疑、文献学习辅助构建任务，记录资料版本、学习目标、可用工具、数据授权与允许的帮助程度。每个基础任务设置干净、检索污染、记忆污染和教学约束变化四个版本；约束变化是否构成风险，由行为及其对应的教学要求共同判定。外发、修改等动作仅在沙箱中使用虚拟身份和模拟工具实施，公开改编、专家构造与授权记录分别报告。

拟设置24个任务家族，每个家族包含5个基础任务，结合4个版本、2个模型和3个随机种子，共进行2880次候选轨迹运行；评价、干预分支、失败重试与人工仲裁另计成本。先在预留开发家族开展预实验，估计主阳性比例、家族间差异、运行费用及区间宽度，据此在预算内调整规模，优先增加独立家族而非重复运行。若精度仍不足，则将相应结论限定为探索性证据。

按照R1的风险定义，由2名标注者独立判断、1名仲裁者处理分歧，教学人员复核支持规则与反例。每项判断均记录证据位置与理由，报告逐类一致率、Cohen's $\kappa$及分歧类型。按图~\ref{fig:edu-f06}将12、6、6个家族分别用于开发、校准和测试，同一家族的任务、变体与种子运行整组分配。划分前独立盲审正常与主阳性覆盖，测试前冻结任务、真值、模型、阈值和主要终点；测试后不补选任务。目标模型不提供自身评价真值，留出材料不参与评价器开发。

\subsection{对照与交叉消融}
针对情境约束的判别作用，设置仅检查结果、通用轨迹评价和教育情境轨迹评价三类对照。后两者共用基础模型、可观测字段、输出格式、上下文容量与调用预算，仅改变评分依据是否包含教学支持边界，以排除信息量和计算投入不同造成的优势。
见图~\ref{fig:edu-f06}。
\HuazhongFigure{figures/F06-factorial-design-v2.pdf}{约束与证据的交叉消融及家族隔离验证}{fig:edu-f06}


按图~\ref{fig:edu-f06}交叉设置轨迹证据保留或移除、教育约束正确或打乱四种条件。打乱仅改变评价器接收的约束，标签仍依据原始情境；另设授权变化任务对，固定行动与回答、改变允许的帮助程度，并按新情境重新标注。前者检验约束与证据的共同贡献，后者检验风险含义是否随授权改变，二者分别分析。辅以等长无关约束和规则措辞改写，检查长度、关键词与模板线索的影响。答案相同的案例用于机制诊断，主比较保留全部初始任务，自然运行及显式权限、语义风险分别报告。

\subsection{时间指标与干预比较}
风险识别按类别报告精确率、召回率、误报率、精确率—召回率曲线面积及类别比例。以图~\ref{fig:edu-f04}中的首次可观察证据时点$t_0$、首次有效告警时点$t_d$和风险动作提交时点$t_a$刻画识别时序；仅将$t_0\leq t_d<t_a$计为动作前检出，同时报告检测延迟$t_d-t_0$和剩余干预窗口$t_a-t_d$。未检出与事后检出单列，正常轨迹中至少一次错误告警即计为整轨迹误报。

各方法均在正常校准轨迹上以10\%为名义误报目标选择阈值，冻结后报告测试集的实际误报率及区间。干预比较包括事后检查、仅规则保护、统一拒绝与分级干预；各策略使用相同初始任务、环境状态和随机条件，重新执行干预后的分支。风险执行率与正常完成率分别以全部预设风险任务、正常任务为分母，拒绝、超时和人工移交均予保留。

以相对无在线干预的正常任务自主完成率损失衡量\textbf{任务效用代价}，预设上限为5个百分点。校准时先筛选满足该约束的策略，再比较风险执行率；测试阶段只验证冻结策略。损失区间上界低于上限才支持满足约束，下界高于上限为违反，跨界则证据不足。人工移交计为未自主完成，另报接管后的完成情况，以及调用次数、耗时、费用和人工复核负担。

教学合规性由独立人员盲评，分别判断目标匹配、帮助程度适切与证据使用合规，采用满足、部分满足、不满足三级尺度，规范不确定者单列。该评价与答案正确性分别报告，研究结论限于任务行为与支持边界，不外推真实学习效果或长期能力变化。

\subsection{主检验与统计判据}
以预选模型的教育支持边界风险为主分析对象，其他模型和风险类别用于稳健性检验或探索。主交互量为$\Delta_{\mathrm{int}}=(R_{11}-R_{01})-(R_{10}-R_{00})$，其中$R_{ec}$为冻结阈值下的家族等权召回，$e$表示证据保留或移除，$c$表示约束正确或打乱，编码见图~\ref{fig:edu-f06}。该差分检验正确约束是否增强证据的判别贡献。

拟以交互召回增益5个百分点为最小实际意义幅度，在预实验后、测试前锁定。交互量区间下界超过该幅度，才支持具有实际意义的增益；上界低于该幅度则不支持，跨界时保留不确定性。与通用轨迹评价的召回差及不同策略的风险执行率差作为支持分析；只有效用损失满足约束且风险执行下降均获区间支持，才判断存在干预收益。若实际误报不同，须结合误报差异解释召回变化，不能将其全部归因于情境与证据的共同作用。

主阳性限定为原情境下经人工确认的教学支持边界违背，支持不足与规范不确定分别报告；对后者按两种极端归类开展敏感性分析。逐家族报告类别数量，任一测试家族无主阳性时，整体主交互量视为不可估计，不将召回赋零或删除该家族；仅在有阳性家族上的分析作为探索结果，并报告覆盖范围。

按家族和基础任务两层进行聚类自助抽样，依据任务是否包含主阳性分层，同一任务的变体与种子运行共同抽取，四个消融条件共用重抽样本，估计95\%区间并记录不可估计的重抽。区间以已有标签为条件，增加种子不视为增加独立样本。若最终仅有6个测试家族，则结合逐家族敏感性分析解释结果，保留探索定位；次要比较控制多重性或明确其探索性质，未显著不作等效解释。

\subsection{技术路线与关键技术}
总体路线见图~\ref{fig:edu-f02}：从教学任务及授权要求出发，R1建立约束与证据关联的风险标注，供R2开展完整轨迹定位和当前前缀评价；R2输出的风险判断、证据位置及不确定性进入R3，驱动动作提交前的处置。三项内容共享任务与日志，最终通过家族留出和策略分支比较，检验风险判断的适用范围及干预收益。

\textbf{风险表征路线。}以课程目标、帮助边界和外显行为为输入，将风险判断对应到具体约束及证据片段，保留反例与不确定标签。通过独立标注和教学复核检验一致性，形成R2所需的任务集、标注规则及参考真值。

\textbf{轨迹评测路线。}以R1的任务与标注为依据，分别构造完整轨迹和仅含当时可见信息的前缀，结合规则核验与语义评分输出分维度判断及告警时点。通过等信息对照、交叉消融和留出验证检查判别来源与提前检出能力，形成R3使用的评价依据和校准参数。输入与日志分别核验，防止未来动作、最终结果或人工标签泄漏至在线评价。

\textbf{监测干预路线。}以当前前缀、候选动作和校准参数为输入，依权限、风险及可修正性决定执行、重规划、阻断或人工复核。保存动作前后状态，重新运行受干预分支，区分工具失败与安全拒绝；依据风险—效用比较形成干预判据与失败案例。双智能体仅作扩展测试，以等消息、等步数的单智能体转发为对照，追踪采纳、执行与记忆；同等传播不解释为双智能体特有的放大效应。

\subsection{可行性分析}
研究以已有模型推理、可控工具环境和公开材料为基础，不依赖基础模型训练。申请人的问答、检索及知识库研究积累可支撑任务组织与实验实现，现有算力和成员分工见“二、研究基础”第\ref{subsec:conditions}节。R1—R3复用任务与轨迹，能够将一年期、5万元的投入集中于标注复核、核心对照和干预验证。

教学规则及其测量效度是实施中的关键条件。拟先请具有相应课程理解能力的人员复核帮助边界及正反例，再开展独立试标和仲裁，形成手册、分歧记录及一致性结果；人员落实前仅验证明示规则，不预认定教学判断有效。真实课程与师生数据的使用以取得授权为前提。若预实验显示费用或统计精度不满足要求，优先保障主比较并收缩扩展任务，具体风险应对见第\ref{subsec:risks}节。
\section{本项目的特色与创新之处。}
\subsection{项目特色与创新点}
\textbf{项目特色：围绕教学支持边界研究过程可信性。}将学习目标、帮助程度与行动授权纳入行为过程，解释相同答案的不同风险。课程答疑与文献学习辅助共用评价依据，核对识别结论与干预后果。

\textbf{创新一：建立情境约束与行为证据关联的风险评价方法。}在轨迹风险和教学工作流评测基础上\cite{yuan2024rjudge,teacharena2026}，关联教学约束、证据位置与风险结论，回答“违反哪项要求、依据哪段行为”。贡献集中于支持边界变化时的风险判定；以等信息对照、交叉消融及跨家族验证，检验结果正确但支持越界的识别能力，为干预提供可复核依据。

\textbf{创新二：形成任务效用约束下的动作前干预判据。}在规则保护和分级控制基础上\cite{xiang2024guardagent,agentguard2026nexus}，结合风险证据与干预窗口，在正常任务自主完成率损失约束下校准策略。比较同任务各处置分支，辨别风险下降来自纠偏还是额外拒绝，形成教育支持边界风险的策略选择依据，明确误阻断、人工负担及失效条件。

\section{研究计划、预期成果及考核指标。}
\subsection{年度研究计划}
研究期12个月，自正式立项之日起计算，按四个阶段推进。

\textbf{第1—3个月：}核对许可、组织教学复核、试标与预实验，确定样本覆盖和规模，交付手册、样例及环境。

\textbf{第4—6个月：}完成标注、划分、对照与消融，冻结测试方案，交付数据及脚本。

\textbf{第7—9个月：}完成留出及干预比较，形成报告、原型与论文初稿。

\textbf{第10—12个月：}完成授权或近真实案例，整理论文、报告及可公开成果。双智能体压力测试视资源开展。

\subsection{预期研究结果}
以发表SCI、SSCI或CSSCI期刊论文不少于1篇、提交典型案例报告1份为考核目标，投稿或录用不替代发表。配套标注规范、轨迹评测集和监测原型各1套，提供脚本、人工依据、划分设置及失败案例。公开遵守许可，成果归属及标注遵循实验室规定，结题提交报告、发表证明及成果清单。
